\documentclass[../../../include/open-logic-section]{subfiles}

\begin{document}

\olfileid{sth}{story}{approach}
\olsection{సంచిత-దశలవారీ పద్ధతి}

సమితులను దశలుగా ఎలా
అమర్చుతామో మరికొంచెం
వివరంగా ఇలా చెప్పవచ్చు:
\begin{quote}
	సమితులు \emph{దశలలో}
	ఏర్పడతాయి. ప్రతి దశ $S$కు,
	$S$కన్నా \emph{ముందు} కొన్ని
	దశలు ఉంటాయి. దశ $S$ వద్ద,
	$S$కన్నా ముందు దశలలో
	ఏర్పడిన సమితులతో కూడిన
	ప్రతి సముదాయం ఒక సమితిగా
	ఏర్పడుతుంది. దశలలో
	ఏర్పడే సమితులు తప్ప
	ఇతర సమితులు లేవు.
	\citep[p.~323]{Shoenfield:AST}
\end{quote} 
ఇది \emph{సమితి సంచిత-దశలవారీ
భావన} యొక్క రూపురేఖ.
\olref[sth][][]{part}లో మనం
చెప్పబోయే నియత సమితి
సిద్ధాంతానికి ఇదే ఆధారం.

దీనిని మరింత వివరంగా
పరిశీలిద్దాం. షోన్‌ఫీల్డ్
వివరించినట్లుగా, ప్రతి దశలో
మునుపటి దశల నుంచి అందుబాటులో
ఉన్న {సమితుల}తో కొత్త సమితులను
ఏర్పరుస్తాం. కాబట్టి ఆయన
చిత్రంలో ప్రారంభ దశ, దశ $0$
వద్ద, \emph{మునుపటి} దశలు
లేవు; అందువల్ల మునుపటి
దశల నుంచి అందుబాటులో
ఉన్న సమితులూ అసలే
లేవు.\footnote{\emph{మొదటి}
దశ ఉందని ఎందుకు అనుకోవాలి?
\olref[z][story]{sec}లోని
\stagesord{} పాదగమనిక
చూడండి.} కాబట్టి ఒకే
సమితిని ఏర్పరుస్తాం:
\tetoken{మూలకాలు}{!!{element}s}
లేని సమితి $\emptyset$.
దశ $1$ వద్ద మునుపటి
దశల నుంచి ఖచ్చితంగా
ఒక సమితి అందుబాటులో
ఉంటుంది; అందువల్ల కొత్త
సమితి ఒక్కటే:
$\{\emptyset\}$.
దశ $2$ వద్ద మునుపటి
దశల నుంచి రెండు సమితులు
అందుబాటులో ఉంటాయి;
రెండు కొత్త సమితులు
$\{\{\emptyset\}\}$,
$\{\emptyset, \{\emptyset\}\}$
ఏర్పడతాయి. దశ $3$
వద్ద నాలుగు సమితులు
అందుబాటులో ఉంటాయి;
అందువల్ల పన్నెండు
కొత్త సమితులు
ఏర్పడతాయి\ldots.
దీనివల్ల సమితుల
సంచిత-దశలవారీ చిత్రం
ఇలా ఉంటుంది (సంఖ్యలు
దశలను సూచిస్తాయి):
\begin{center}
	\begin{tikzpicture}[scale=0.6]
	\tikzset{cut_here/.style={densely dotted}}
	\draw (-4,6) -- (-1, 0)-- (2,6); 
	\draw[cut_here] (-5,8)--(-4,6);
	\draw[cut_here] (2,6)--(3,8);
	\draw(-1.5, 1)--(-0.5, 1);
	\draw(-2, 2)--(0, 2);
	\draw(-2.5, 3)--(0.5,3);
	\draw(-3, 4)--(1, 4);
	\draw(-3.5, 5)--(1.5, 5);
	\draw(-4, 6)--(2, 6);
	\node[label] at (0, 0) {\small 0};
	\node[label] at (0.5, 1) {\small 1};
	\node[label] at (1, 2) {\small 2};
	\node[label] at (1.5, 3) {\small 3};
	\node[label] at (2, 4) {\small 4};
	\node[label] at (2.5, 5) {\small 5};
	\node[label] at (3, 6) {\small 6};
	\end{tikzpicture}
\end{center}
అయితే ఈ కథనాన్ని
ఎందుకు స్వీకరించాలి?

ఒక కారణం: ఇది సరళంగా
అర్థమయ్యే, నిర్వహించదగిన
కథనం. అత్యంత స్పష్టంగా
అనిపించిన కథనం,
అంటే నిర్బంధరహిత
ధర్మసంగ్రహం, విఫలమైన
తరువాత మనకు మంచి
ప్రత్యామ్నాయం కావాలి.

అయితే ఈ కథనం
\emph{కేవలం} ఆకర్షణీయమైనది
మాత్రమే కాదు. దీనిపై
ఆధారపడిన సమితి సిద్ధాంతం
\emph{అవైరుధ్యమైనది} అని
నమ్మడానికి మంచి కారణం
ఉంది. అదేమిటంటే:

సమితి సంచిత-దశలవారీ
భావన ప్రకారం సమితులు
దశల్లో ఏర్పడతాయి;
వాటి \tetoken{మూలకాలు}{!!{element}s}
\emph{ఇప్పటికే} అందుబాటులో
ఉన్న వస్తువులై ఉండాలి.
కాబట్టి ఏ దశ~$S$కైనా
కింది సమితిని ఏర్పరచవచ్చు:
\[
	R_S = \Setabs{x}{x \notin x \text{ మరియు $x$ అనేది $S$కు ముందు అందుబాటులో ఉంది}}
\]
రసెల్ వైరుధ్యాన్ని
నిరూపించే వాదన ఇప్పుడు
$R_S$ స్వయంగా దశ
$S$కు ముందు అందుబాటులో
లేదని స్థాపిస్తుంది.
అది వైరుధ్యం కాదు.
అంతేకాకుండా, సమితి
సంచిత-దశలవారీ భావనను
అంగీకరిస్తే రసెల్
సమితినే ఏర్పరచగలమని
\emph{ఆశించకూడదు}.
ఎందుకంటే అది ``ఎప్పుడైనా
అందుబాటులోకి వచ్చే''
తమలో తాము మూలకాలుగా
లేని అన్ని సమితుల
సమితి అవుతుంది.
సంక్షిప్తంగా, రసెల్
సమితిని (నిరూపితంగా)
ఏర్పరచలేమన్న విషయం
సంచిత-దశలవారీ కథనంలో
\emph{ఆశ్చర్యకరం కాదు};
మనం \emph{ఊహించేదే}.

%ఒక అర్థంలో, సమితి సంచిత-దశలవారీ భావన రసెల్ వైరుధ్యానికి ఇచ్చే సమాధానం \emph{స్వవర్గానవలంబనవాది} సమాధానంలాంటిదే. తమలో తాము మూలకాలుగా లేని సమితులు$_n$ అన్నింటి సముదాయం సమితి$_{n+1}$ అని, ఆ సమితి$_n$ తనలో తాను మూలకంగా ఉండదని ఆయన అంటారు. (నిజానికి, సమితి తనలో తాను మూలకమా అని అడగడమే \emph{వ్యాకరణవిరుద్ధం} అని చాలామంది అలాంటి వాదులు భావిస్తారు.)
%
%అయితే సంచిత-దశలవారీ పద్ధతికీ స్వవర్గానవలంబనవాది పద్ధతికీ ముఖ్య భేదం ఉంది: మొదటిది అన్ని వస్తువులనూ \emph{ఒకే రకానికి చెందినవి}గా పరిగణిస్తుంది. రెండో పద్ధతిలో శూన్య సమితి$_0$ (అంటే \tetoken{మూలకాలు}{!!{element}s} లేని సమితి$_0$), శూన్య సమితి$_1$ (అంటే \tetoken{మూలకాలు}{!!{element}s} లేని సమితి$_1$) ఉంటాయి; అవి \emph{వేర్వేరు వస్తువులు}. కానీ సంచిత-దశలవారీ పద్ధతిలో ఒకే శూన్య సమితి $\emptyset$ ఉంటుంది; అది ప్రతి దశలో అందుబాటులో ఉంటుంది.

%నిజానికి ఈ అధ్యాయం మొదట చెప్పిన సమితుల సమానత్వ స్వీకృతం ఈ శ్రేణిలోని మూలకాలకు వర్తిస్తుంది (కాబట్టి \emph{ఒక్కటే} ``శూన్య'' సమితి ఉంటుంది). అయితే సంచిత-దశలవారీ భావనతో ఆ స్వీకృతాన్ని సమర్థించాలనుకోవడం లేదు (\cite{Boolos1971} చూడండి). మళ్లీ చెప్పాలంటే, మన ఆసక్తి మూలకాధారిత సముదాయాలపై కాబట్టే సమితుల సమానత్వాన్ని అంగీకరించాలి.

\end{document}
